\documentclass[12pt]{article}
\usepackage{amsmath,amssymb}
\usepackage[margin=1in]{geometry}

\begin{document}
	
	\section*{Annotated explanation of the definition of a full-rank lattice}
	
	We are given the statement:
	
	\[
	\text{A (full-rank) lattice } \Lambda \subset \mathbb{R}^n \text{ is }
	\Lambda = \{B\mathbf{z} : \mathbf{z}\in\mathbb{Z}^n\}
	\]
	for some basis matrix \(B\in\mathbb{R}^{n\times n}\).
	
	\bigskip
	
	This means the following.
	
	\subsection*{1. What is \(B\)?}
	
	The matrix \(B\in\mathbb{R}^{n\times n}\) is an \(n\times n\) real matrix whose columns are
	\[
	B=
	\begin{bmatrix}
		\mathbf{b}_1 & \mathbf{b}_2 & \cdots & \mathbf{b}_n
	\end{bmatrix}.
	\]
	
	So the columns of \(B\) are vectors
	\[
	\mathbf{b}_1,\mathbf{b}_2,\dots,\mathbf{b}_n\in\mathbb{R}^n.
	\]
	
	Since \(B\) is called a \emph{basis matrix}, its columns are linearly independent. Hence they form a basis of \(\mathbb{R}^n\).
	
	\subsection*{2. What is \(\mathbf{z}\in\mathbb{Z}^n\)?}
	
	The vector \(\mathbf{z}\) is an integer vector:
	\[
	\mathbf{z}=
	\begin{pmatrix}
		z_1\\
		z_2\\
		\vdots\\
		z_n
	\end{pmatrix},
	\qquad z_1,z_2,\dots,z_n\in\mathbb{Z}.
	\]
	
	So each coordinate of \(\mathbf{z}\) must be an integer.
	
	\subsection*{3. What does \(B\mathbf{z}\) mean?}
	
	When we multiply \(B\) by \(\mathbf{z}\), we get
	\[
	B\mathbf{z}
	=
	z_1\mathbf{b}_1+z_2\mathbf{b}_2+\cdots+z_n\mathbf{b}_n.
	\]
	
	So \(B\mathbf{z}\) is an \emph{integer linear combination} of the columns of \(B\).
	
	\subsection*{4. What is the set \(\Lambda=\{B\mathbf{z}:\mathbf{z}\in\mathbb{Z}^n\}\)?}
	
	This means:
	\[
	\Lambda
	=
	\left\{
	z_1\mathbf{b}_1+z_2\mathbf{b}_2+\cdots+z_n\mathbf{b}_n
	:\ z_1,\dots,z_n\in\mathbb{Z}
	\right\}.
	\]
	
	So \(\Lambda\) is the set of all points you can get by taking integer combinations of the basis vectors \(\mathbf{b}_1,\dots,\mathbf{b}_n\).
	
	\subsection*{5. Why is it called a lattice?}
	
	Because the resulting set is a discrete, regularly repeated grid of points in \(\mathbb{R}^n\).
	
	It is like the usual integer grid \(\mathbb{Z}^n\), except it may be stretched, rotated, or skewed by the matrix \(B\).
	
	\subsection*{6. Why ``full-rank''?}
	
	The phrase \emph{full-rank} means that \(B\) has rank \(n\). Equivalently, its columns are linearly independent.
	
	Therefore the lattice is not trapped in a lower-dimensional subspace. It lives in all \(n\) dimensions of \(\mathbb{R}^n\).
	
	\subsection*{7. Compare with the real span}
	
	If we allowed real coefficients instead of integer coefficients, we would get
	\[
	\{B\mathbf{a}:\mathbf{a}\in\mathbb{R}^n\}
	=
	\mathbb{R}^n
	\]
	because \(B\) is invertible.
	
	But in the lattice definition we only allow integer vectors \(\mathbf{z}\in\mathbb{Z}^n\), so we get only a discrete subset:
	\[
	\Lambda=\{B\mathbf{z}:\mathbf{z}\in\mathbb{Z}^n\}.
	\]
	
	Thus:
	\begin{itemize}
		\item real coefficients give a whole continuous space,
		\item integer coefficients give a discrete grid of lattice points.
	\end{itemize}
	
	\subsection*{8. Standard example}
	
	If
	\[
	B=I_n,
	\]
	the identity matrix, then
	\[
	\Lambda=\{I_n\mathbf{z}:\mathbf{z}\in\mathbb{Z}^n\}=\mathbb{Z}^n.
	\]
	
	So the standard integer lattice is the simplest example of a full-rank lattice.
	
	\subsection*{9. Example in \(\mathbb{R}^2\)}
	
	Let
	\[
	B=
	\begin{pmatrix}
		2 & 1\\
		0 & 1
	\end{pmatrix}.
	\]
	Then its columns are
	\[
	\mathbf{b}_1=
	\begin{pmatrix}
		2\\
		0
	\end{pmatrix},
	\qquad
	\mathbf{b}_2=
	\begin{pmatrix}
		1\\
		1
	\end{pmatrix}.
	\]
	
	The lattice is
	\[
	\Lambda
	=
	\{B\mathbf{z}:\mathbf{z}\in\mathbb{Z}^2\}.
	\]
	
	If
	\[
	\mathbf{z}=
	\begin{pmatrix}
		z_1\\
		z_2
	\end{pmatrix},
	\qquad z_1,z_2\in\mathbb{Z},
	\]
	then
	\[
	B\mathbf{z}
	=
	\begin{pmatrix}
		2 & 1\\
		0 & 1
	\end{pmatrix}
	\begin{pmatrix}
		z_1\\
		z_2
	\end{pmatrix}
	=
	\begin{pmatrix}
		2z_1+z_2\\
		z_2
	\end{pmatrix}.
	\]
	
	Therefore
	\[
	\Lambda
	=
	\left\{
	\begin{pmatrix}
		2z_1+z_2\\
		z_2
	\end{pmatrix}
	: z_1,z_2\in\mathbb{Z}
	\right\}.
	\]
	
	This is a skewed infinite grid in the plane.
	
	\subsection*{10. Plain English summary}
	
	The definition
	\[
	\Lambda=\{B\mathbf{z}:\mathbf{z}\in\mathbb{Z}^n\}
	\]
	means:
	
	\begin{quote}
		Choose \(n\) linearly independent vectors in \(\mathbb{R}^n\).  
		Then take all possible whole-number combinations of those vectors.  
		The set of resulting points is the lattice.
	\end{quote}
	
\end{document}